% ============================================================
%  sections/07_conclusion.tex
% ============================================================
\section{Conclusion}
\label{sec:conclusion}

We presented a controlled empirical benchmark of three neural network
compression methods—magnitude pruning, dynamic INT8 quantization, and
SVD low-rank factorization—on a compact CNN trained on Fashion-MNIST.
Our contributions are threefold.

\textbf{Theoretically}, we proved three results:
sparsity preservation under masked gradient descent (Theorem~\ref{thm:sparsity}),
the Eckart-Young optimality of truncated SVD (Theorem~\ref{thm:eckart_young}),
and INT8 quantization error and network perturbation bounds (Theorem~\ref{thm:quant_error}).
These results give a principled foundation for understanding why each
method behaves as it does.

\textbf{Empirically}, we uncovered a \textbf{non-monotonic latency profile}
for unstructured pruning that contradicts the simplified ``no speedup'' narrative:
moderate sparsity (25--75\%) \emph{increases} wall-clock latency by 13--41\%
due to masking overhead and cache-access degradation, while extreme sparsity
(90\%) produces an unexpected 40\% latency \emph{reduction} explained by
L2 cache residency effects. Neither trend is predictable from theoretical
FLOPs analysis alone—a finding with direct implications for how compression
benchmarks should be reported.

\textbf{Practically}, our Pareto analysis identifies dynamic INT8
quantization as the dominant compression strategy on commodity CPU
hardware: it achieves $\sim$4$\times$ model size reduction with
$<$0.1\% accuracy degradation and is the only method that reduces
\emph{both} model size and inference latency simultaneously.

\paragraph{Future Work.}
Promising extensions include:
structured channel pruning for real latency benefits;
quantization-aware training (QAT) for improved accuracy under
aggressive quantization;
architectural low-rank factorization;
and combined compression pipelines.

\paragraph{Code Availability.}
All code, experiment configurations, and benchmark data are available at:
\url{https://github.com/your-username/AI_Optimise}
